\documentclass{article}

% Self-published preprint. This paper was not accepted at (or submitted to)
% ICML 2026; the ICML style file is used for formatting only.
\usepackage[preprint]{icml2026}
\usepackage{amsmath,amssymb,amsthm}
\usepackage{graphicx}
\usepackage{booktabs}
\usepackage{algorithm}
\usepackage{algorithmic}
\usepackage{hyperref}
\usepackage{xcolor}
\usepackage{subcaption}

\newtheorem{theorem}{Theorem}
\newtheorem{definition}{Definition}
\newtheorem{proposition}{Proposition}
\newtheorem{lemma}{Lemma}
\newtheorem{corollary}{Corollary}

\icmltitlerunning{Hodge Potential Alignment for Preference Optimization: A Negative Result}

\begin{document}

\twocolumn[
\icmltitle{Hodge Potential Alignment for Preference Optimization:\\A Negative Result and Three Evaluation Pitfalls}

\begin{icmlauthorlist}
\icmlauthor{Mike H Lee}{ox}
\end{icmlauthorlist}

\icmlaffiliation{ox}{O-X}

\icmlcorrespondingauthor{Mike H Lee}{mike@oasis-x.io}

\icmlkeywords{preference optimization, Hodge decomposition, cyclic preferences, negative result, evaluation}

\vskip 0.3in
]

\printAffiliationsAndNotice{}

\begin{abstract}
Pairwise preference data can contain cycles ($A \succ B \succ C \succ A$) that no scalar reward can represent. We tested \textbf{HodgePO}, a method that computes the combinatorial Hodge decomposition of a preference graph and adds a regularizer that pulls a model's reward differences toward the cycle-free (gradient) potential. An earlier version of this work reported large gains: in-sample ranking accuracy rose from $0.940$ to $0.9999$ for DPO and from $0.800$ to $0.9964$ for KTO over 30 seeds. \textbf{Those gains do not survive a controlled re-test.} On held-out pairs (5 train/held-out splits $\times$ 30 seeds, 500 HH-RLHF pairs), Hodge-DPO scores $0.537$ and DPO $0.550$; Hodge-KTO $0.543$ and KTO $0.542$. No paired difference is significant at the split level ($p \geq 0.34$), and every method, including a linear probe, stays between $0.52$ and $0.59$. On the training pairs, a \emph{margin control}---the same loss term with every target set to the mean Hodge target---reproduces the whole in-sample gain. We also found that the original run gave each training pair the Hodge target of a different pair, and that the gain did not depend on it. We report the negative result together with the three evaluation pitfalls that produced the original claim, so that others who apply Hodge methods to preference data can avoid them.
\end{abstract}

%% ====================================================================
\section{Introduction}
\label{sec:intro}
%% ====================================================================

Preference optimization methods---Direct Preference Optimization (DPO)~\citep{rafailov2023direct}, Kahneman-Tversky Optimization (KTO)~\citep{ethayarajh2024kto}, Group Relative Policy Optimization (GRPO)~\citep{shao2024deepseekmath}, and Odds Ratio Preference Optimization (ORPO)~\citep{hong2024orpo}---are widely used to align language models with pairwise preference data. Most of them assume, explicitly or implicitly, a scalar reward.

Pairwise preferences are not always consistent with any scalar reward. Majority preferences can be cyclic even when every individual preference is transitive~\citep{condorcet1785essai, arrow1950difficulty}. \citet{liu2025impossibility} prove that preferences among LLM responses are representable by a reward model if and only if they contain no Condorcet cycle. The combinatorial Hodge decomposition~\citep{hodge1941harmonic, jiang2011statistical} splits a pairwise flow into a part that a scalar potential explains exactly and a cyclic remainder. A natural idea is to use the potential as a training target.

We tested that idea and first reported a strong positive result. A controlled re-test shows that the result was an artifact of evaluation. This paper reports:
\begin{enumerate}
    \item \textbf{The method} (Section~\ref{sec:method}): per-sample weights from cycle participation, and a potential-alignment regularizer.
    \item \textbf{The negative result} (Section~\ref{sec:heldout}): on held-out pairs, HodgePO does not beat its base optimizers or a margin control that carries no pair-specific Hodge information.
    \item \textbf{Three evaluation pitfalls} (Section~\ref{sec:pitfalls}) that together produced the original claim: in-sample scoring, a target that acts as a generic margin, and an index misalignment between samples and graph edges.
    \item \textbf{Two practical facts} about Hodge methods on preference data: a graph with one edge per pair has no cycles, and a Hodge penalty on the outputs of a scalar reward model is identically zero (Appendix~\ref{app:batch_penalty}).
\end{enumerate}

%% ====================================================================
\section{Background}
\label{sec:background}
%% ====================================================================

\subsection{Preference Optimization}

Given preference pairs $\mathcal{D} = \{(x, y_w, y_l)\}$ with $y_w \succ y_l$ for prompt $x$, DPO~\citep{rafailov2023direct} minimizes
\begin{equation}
    \mathcal{L}_\text{DPO} = -\mathbb{E}\left[\log \sigma\left(\beta \log \frac{\pi_\theta(y_w | x)}{\pi_\text{ref}(y_w | x)} - \beta \log \frac{\pi_\theta(y_l | x)}{\pi_\text{ref}(y_l | x)}\right)\right].
\label{eq:dpo}
\end{equation}
KTO~\citep{ethayarajh2024kto} scores individual samples with prospect-theoretic weighting; GRPO~\citep{shao2024deepseekmath} uses group-relative advantages. These methods are consistent with a Bradley-Terry model~\citep{bradley1952rank}, $P(y_w \succ y_l) = \sigma(r^*(y_w) - r^*(y_l))$, which cannot represent cycles.

\subsection{Combinatorial Hodge Theory}
\label{sec:hodge_background}

Let $G = (V, E)$ be a graph with an orientation on each edge and let $K$ be its clique complex. An edge flow is a 1-cochain $f \in C^1$. Let $d_0: C^0 \to C^1$, $(d_0 \phi)_{ij} = \phi(j) - \phi(i)$, and $d_1: C^1 \to C^2$, $(d_1 f)_{ijk} = f_{ij} + f_{jk} + f_{ki}$.

\begin{theorem}[Hodge decomposition~\citep{jiang2011statistical}]
\label{thm:hodge}
Every edge flow decomposes uniquely and orthogonally as $f = d_0 \phi + d_1^{*} \psi + h$, with $\phi \in C^0$, $\psi \in C^2$, and $h \in \ker \Delta_1$, where $\Delta_1 = d_0 d_0^{*} + d_1^{*} d_1$.
\end{theorem}

The gradient part $d_0\phi$ is what a global ranking explains; the curl and harmonic parts form the cyclic remainder. The gradient potential is the least-squares solution
\begin{equation}
    \phi^* = \arg\min_\phi \sum_{(i,j) \in E} \left(f_{ij} - (\phi(j) - \phi(i))\right)^2 .
\end{equation}
Our implementation adds no 2-cells, so it computes $\phi^*$ and the cyclic remainder $r = f - d_0\phi^*$ (curl and harmonic together). Each edge's flow is the log-odds of its preference probability, clipped to $[-10, 10]$.

%% ====================================================================
\section{Method}
\label{sec:method}
%% ====================================================================

\subsection{Preference Graph}
\label{sec:graph_construction}

With one direct edge per training pair, and each response appearing in one pair only, the graph is a matching. A matching has no cycles, so the decomposition is trivial. All cyclic structure must come from edges added between pairs.

\begin{definition}[Cross-pair preference graph]
Given $N$ pairs with response embeddings $\{(\mathbf{e}_w^{(i)}, \mathbf{e}_l^{(i)})\}_{i=1}^N$, the graph has:
\begin{itemize}
    \item \textbf{Direct edges} $(w_i, l_i)$, with the pair's annotation confidence as preference probability.
    \item \textbf{Similarity edges}: for each preferred response $w_i$, its $k = 5$ nearest preferred responses $w_j$ (cosine similarity $s_{ij}$) from other pairs; an edge $w_i \to w_j$ with probability $\mathrm{clip}(0.5 + 0.2(s_{ij} - 0.5), 0.1, 0.9)$, and an edge $w_i \to l_j$ with probability $\mathrm{clip}(0.7 + 0.2(1 - s'_{ij}), 0.55, 0.95)$, where $s'_{ij}$ is the similarity of $w_i$ and $l_j$.
\end{itemize}
\end{definition}

The similarity-edge probabilities are fixed formulas, not human judgments. The cycles that the decomposition finds are therefore a property of this construction, and they say nothing about how cyclic human preferences are.

\subsection{Weights and Regularizer}
\label{sec:potential_alignment}

For node $v$, let $s(v) = \sum_{e \ni v} |r_e|$. For pair $i$, $\text{cycle\_score}(i) = (s(w_i) + s(l_i)) / \max_j (s(w_j) + s(l_j))$, and the weight is $\alpha_i = 1 - 0.6 \cdot \text{cycle\_score}(i)$. The potential target is
\begin{equation}
    \Delta\phi_i = \phi^*(l_i) - \phi^*(w_i),
\label{eq:potential_diff}
\end{equation}
which is positive when the pair agrees with the global ranking (flow on $w_i \to l_i$ encodes $w_i \succ l_i$). The regularizer is
\begin{equation}
    \mathcal{L}_\text{Hodge} = \frac{1}{B} \sum_{i} \left(r_\theta(w_i) - r_\theta(l_i) - \Delta\phi_i\right)^2 ,
\label{eq:hodge_loss}
\end{equation}
and the total loss is $\mathcal{L}_\text{base}(\alpha) + \lambda_\text{Hodge} \mathcal{L}_\text{Hodge}$ with $\lambda_\text{Hodge} = 0.05$.

%% ====================================================================
\section{Experiments}
\label{sec:experiments}
%% ====================================================================

\subsection{Common Setup}

All methods train a 2-layer MLP ($384 \to 64 \to 1$) over frozen \texttt{all-MiniLM-L6-v2} embeddings of prompts and responses, with learning rate $10^{-3}$, batch size 64, and 50 epochs. The metric is \textbf{ranking accuracy}: the fraction of pairs on which the model gives the preferred response the higher reward. (Our code names this quantity ``exploit resistance''. It does not measure reward hacking: there is no environment, no text generation, and no policy optimized against the learned reward.)

\subsection{The Original Run (In-Sample)}
\label{sec:original}

The original run sampled 500 pairs from a pool of 2,268 pairs (2,000 HH-RLHF harmless-base pairs~\citep{bai2022training} and 268 LLM-written counterfactual pairs from the TRACE data set), built the graph from the kept pairs, and trained and scored every method on the same 500 pairs, over 30 seeds. Table~\ref{tab:original} gives the result.

\begin{table}[t]
\centering
\caption{Original run: ranking accuracy on the training pairs (in-sample), 30 seeds. These numbers do not survive the re-test in Section~\ref{sec:heldout}.}
\label{tab:original}
\vskip 0.1in
\begin{tabular}{@{}lcc@{}}
\toprule
\textbf{Method} & \textbf{Mean} & \textbf{Std} \\
\midrule
DPO & $0.9403$ & $0.0129$ \\
Hodge-DPO & $0.9999$ & $0.0005$ \\
KTO & $0.8004$ & $0.0166$ \\
Hodge-KTO & $0.9964$ & $0.0026$ \\
GRPO, Hodge-GRPO & $1.000$ & $0.000$ \\
ORPO & $0.6217$ & $0.0803$ \\
\bottomrule
\end{tabular}
\end{table}

The subsample was drawn from an unseeded random generator, so this run cannot be reproduced exactly. A replay of 200 seeded draws with the same code gives baselines whose range contains the published values (DPO $0.907$--$0.963$; KTO $0.735$--$0.827$; 24 draws $\times$ 30 seeds), so the table is typical of the pipeline.

\subsection{Controlled Re-Test (Held-Out)}
\label{sec:heldout}

We used 500 HH-RLHF harmless-base pairs and made 5 random 80/20 train/held-out splits. For each split, the graph, the nearest-neighbour index, the dimensionality reduction, the Hodge potential, and the weights are computed from the 400 training pairs only; held-out texts pass only through the frozen encoder. Each train graph has 4,398--4,400 edges (400 direct). Every method trains 30 seeds per split. We add three arms:
\begin{itemize}
    \item \textbf{Margin control}: the same regularizer and $\lambda$, but every pair's target is the mean of the Hodge targets over the training pairs, and all weights are 1. It has the same scale as HodgePO and no pair-specific Hodge information.
    \item \textbf{Weights only}: HodgePO with $\lambda_\text{Hodge} = 0$.
    \item \textbf{Misaligned}: HodgePO with targets and weights permuted across pairs (see Section~\ref{sec:pitfalls}).
\end{itemize}

\begin{table}[t]
\centering
\caption{Controlled re-test. Ranking accuracy, mean $\pm$ sd over 150 runs (5 splits $\times$ 30 seeds). Chance is $0.50$. A linear Bradley-Terry probe on the same features scores $0.54$--$0.58$ held-out.}
\label{tab:heldout}
\vskip 0.1in
\resizebox{\columnwidth}{!}{%
\begin{tabular}{@{}lcc@{}}
\toprule
\textbf{Method} & \textbf{Held-out} & \textbf{Train} \\
\midrule
DPO               & $0.550 \pm 0.038$ & $0.922 \pm 0.015$ \\
Hodge-DPO         & $0.537 \pm 0.030$ & $1.000 \pm 0.000$ \\
Margin control (DPO) & $0.544 \pm 0.030$ & $1.000 \pm 0.000$ \\
Weights only (DPO) & $0.546 \pm 0.038$ & $0.922 \pm 0.015$ \\
Misaligned (DPO)  & $0.540 \pm 0.031$ & $1.000 \pm 0.000$ \\
\midrule
KTO               & $0.542 \pm 0.051$ & $0.775 \pm 0.018$ \\
Hodge-KTO         & $0.543 \pm 0.027$ & $0.998 \pm 0.002$ \\
Margin control (KTO) & $0.546 \pm 0.030$ & $0.997 \pm 0.004$ \\
Weights only (KTO) & $0.541 \pm 0.049$ & $0.774 \pm 0.020$ \\
Misaligned (KTO)  & $0.547 \pm 0.028$ & $0.998 \pm 0.003$ \\
\midrule
GRPO              & $0.554 \pm 0.034$ & $1.000 \pm 0.000$ \\
ORPO              & $0.535 \pm 0.041$ & $0.638 \pm 0.077$ \\
\bottomrule
\end{tabular}%
}
\end{table}

\begin{table}[t]
\centering
\caption{Paired held-out differences. The 30 seeds in a split share one held-out set, so the split-level test ($n = 5$ split means) is the fair test.}
\label{tab:paired}
\vskip 0.1in
\resizebox{\columnwidth}{!}{%
\begin{tabular}{@{}lccc@{}}
\toprule
\textbf{Comparison} & \textbf{Mean diff.} & \textbf{Splits $>0$} & \textbf{Split-level $p$} \\
\midrule
Hodge-DPO $-$ DPO & $-0.013$ & 1/5 & $0.49$ \\
Hodge-DPO $-$ margin control & $-0.006$ & 2/5 & $0.34$ \\
Hodge-KTO $-$ KTO & $+0.001$ & 2/5 & $0.97$ \\
Hodge-KTO $-$ margin control & $-0.003$ & 2/5 & $0.52$ \\
\bottomrule
\end{tabular}%
}
\end{table}

\paragraph{Result.} On held-out pairs, HodgePO does not beat its base method or the margin control (Tables~\ref{tab:heldout} and~\ref{tab:paired}). At the seed level, where tests overstate the evidence, Hodge-DPO is slightly \emph{worse} than DPO (mean $-0.013$, $p = 0.002$). On the training pairs, HodgePO raises accuracy on every run (DPO: $+0.078$, 150/150; KTO: $+0.223$, 150/150), and the margin control reproduces this gain almost exactly (Hodge minus margin control: $0.0000$ for DPO, $+0.002$ for KTO). The in-sample gain therefore comes from the extra positive-margin term, not from pair-specific Hodge information.

\paragraph{Power.} Every method, including the linear probe, reaches only $0.52$--$0.59$ held-out, so this benchmark cannot resolve small differences. The data do not rule out a small Hodge effect; they give no evidence for one.

%% ====================================================================
\section{Three Evaluation Pitfalls}
\label{sec:pitfalls}
%% ====================================================================

\paragraph{1. In-sample scoring.} The original benchmark scored each model on the pairs it trained on. With 500 pairs and a model that can memorize them, in-sample accuracy measures fit, not generalization: train accuracy is $0.92$--$1.00$ while held-out accuracy is $0.54$--$0.55$ for the same models.

\paragraph{2. A target that acts as a generic margin.} All Hodge targets are positive (in the re-test, 400 of 400 in every split; split means $0.40$--$0.41$, minimum $0.12$). A squared loss toward a positive target pushes every pair toward a positive reward difference, which raises in-sample ranking accuracy regardless of what the targets encode. Without a margin control, this effect is indistinguishable from a Hodge effect.

\paragraph{3. Index misalignment.} In the original run, the graph's edges kept their original order after subsampling, while the training samples followed the random draw order, and the code read sample $i$'s target from edge $i$. In a replay of the subsampling (200 draws), on average $1.0$ of 500 samples received their own edge's target (maximum 6). The original run therefore trained each pair toward another pair's target. The in-sample gain survived this error unchanged (Misaligned arms, Table~\ref{tab:heldout}), which is itself evidence that the targets carried no useful pair-specific information.

\paragraph{Checks we recommend.} Score on held-out pairs; include a control that has the same loss form without the structural information; and assert, in code, that each sample's structural target belongs to that sample.

%% ====================================================================
\section{Related Work}
\label{sec:related}
%% ====================================================================

\paragraph{Preference optimization beyond scalar rewards.} Nash Learning from Human Feedback~\citep{munos2023nash} and the minimax-winner approach of \citet{swamy2024minimaximalist} replace the scalar reward with a preference model and a game-theoretic solution concept. \citet{zhang2025gpm} represent cyclic preferences explicitly. \citet{liu2025impossibility} prove that a reward model exists if and only if the preferences have no Condorcet cycle.

\paragraph{Closest work.} \citet{huang2026transitivity} decompose preferences into orthogonal transitive and cyclic components for LLM alignment, building on the game decomposition of \citet{balduzzi2019open}; \citet{candogan2011flows} give the potential/harmonic decomposition of games. Our negative result concerns one specific use of the decomposition---the gradient potential as a regression target on an embedding-level benchmark---and does not test these methods.

\paragraph{Hodge theory in ranking and ML.} HodgeRank~\citep{jiang2011statistical} introduced the combinatorial Hodge decomposition for ranking. Sheaf methods appear in graph neural networks~\citep{bodnar2022neural}; \citet{ayzenberg2025sheaf} survey the area.

\paragraph{Reward overoptimization.} \citet{moskovitz2024confronting} and \citet{coste2024reward} address overoptimization of a learned reward during policy optimization, which our benchmark does not include.

%% ====================================================================
\section{Discussion and Limitations}
\label{sec:discussion}
%% ====================================================================

\paragraph{Scope of the negative result.} The re-test uses one data set (HH-RLHF harmless-base), 500 pairs, one frozen encoder, and a small MLP; held-out accuracy for all methods is close to chance. The result shows that this benchmark provides no evidence for HodgePO. It does not show that Hodge-derived targets can never help, for example with learned encoders, larger data, or data with genuinely repeated comparisons, where cycles are observed rather than constructed.

\paragraph{What remains useful.} Two facts hold independently of the benchmark: a graph with one edge per pair has no cycles, and a Hodge penalty on scalar predictions is identically zero. The three pitfalls in Section~\ref{sec:pitfalls} apply to any method that injects precomputed structural targets into training.

%% ====================================================================
\section{Conclusion}
\label{sec:conclusion}
%% ====================================================================

We reported large in-sample gains for a Hodge potential-alignment regularizer, then tested them properly. On held-out pairs the regularizer does not beat its base optimizers or a margin control, and the original gain is fully explained by an extra positive-margin term. We publish the negative result, the replay of the original run, and the three pitfalls that produced it.

\bibliography{references}
\bibliographystyle{icml2026}

%% ====================================================================
\newpage
\appendix

\section{Proof Sketch of Theorem~\ref{thm:hodge}}
\label{app:hodge_proof}

With the inner product $\langle f, g \rangle = \sum_e f_e g_e$ on $C^1$, $d_1 d_0 = 0$ implies that $\operatorname{im} d_0$ and $\operatorname{im} d_1^{*}$ are orthogonal, and the finite-dimensional Hodge theorem gives $C^1 = \operatorname{im}(d_0) \oplus \operatorname{im}(d_1^{*}) \oplus \ker(\Delta_1)$. The potential solves $d_0^{*} d_0 \phi = d_0^{*} f$, unique up to a constant on each connected component; we add a ridge term $10^{-4} I$.

\section{Why a Batch Harmonic Penalty Is Zero}
\label{app:batch_penalty}

For batch rewards $\{r_1, \dots, r_{2B}\}$ and the complete graph with flow $f_{ij} = r_j - r_i$, the potential $\phi_i = r_i$ reproduces every edge: $(d_0\phi)_{ij} = f_{ij}$. The cyclic remainder is zero for all $(i,j)$. A flow induced by scalar predictions is a pure gradient, so the decomposition can find cycles only in flows that come from outside the model.

\section{Hyperparameters}
\label{app:implementation}

\begin{table}[h]
\centering
\begin{tabular}{@{}ll@{}}
\toprule
\textbf{Parameter} & \textbf{Value} \\
\midrule
Embedding (frozen) & all-MiniLM-L6-v2, 384 \\
Hidden dim & 64 \\
Learning rate & $10^{-3}$ \\
Epochs & 50 \\
Batch size & 64 \\
DPO $\beta$, KTO $\beta$ & 0.1, 0.1 \\
GRPO $\beta$, group size & 0.04, 8 \\
ORPO $\lambda$ & 0.5 \\
$\lambda_\text{Hodge}$ & 0.05 \\
Weight damping & 0.6 \\
$k$ (nearest neighbours) & 5 \\
\bottomrule
\end{tabular}
\end{table}

\section{Provenance}
\label{app:provenance}

{\sloppy
All files are in the \texttt{shape\_of\_good\_behavior} repository. Table~\ref{tab:original}: \url{shared/results/optimizer_comparison_hodge_v3_30seed.json}, produced by \texttt{run\_optimizer\_comparison} in \url{shared/modal_runner.py} from a 2,268-pair mapping stored on a Modal volume. Replay of that run: \url{scripts/replay_published_v3_subsample.py}. Tables~\ref{tab:heldout} and~\ref{tab:paired}: \url{shared/results/optimizer_comparison_heldout_v1.json}, produced by \url{scripts/heldout_hodge_benchmark.py} from \url{shared/results/counterfactual_pairs.json}. The file \url{shared/results/README.md} records the full history of this benchmark.\par}

\end{document}
